\documentclass{article}

\usepackage[nonanonymous,final]{neurips_2026}

\usepackage[utf8]{inputenc}
\usepackage[T1]{fontenc}
\usepackage{hyperref}
\usepackage{url}
\usepackage{booktabs}
\usepackage{amsfonts}
\usepackage{nicefrac}
\usepackage{microtype}
\usepackage{xcolor}
\usepackage{graphicx}

\title{Food Image Classification Using Transfer Learning on Food-101}

\author{%
  Jasia Azreen \\
  90739129 \\
  Department of Electrical and Computer Engineering \\
  University of British Columbia \\
  \And
  John Song \\
  12837472 \\
  Department of Electrical and Computer Engineering \\
  University of British Columbia \\
  \And
  Sofiya Spolitak \\
  69202497 \\
  Department of Electrical and Computer Engineering \\
  University of British Columbia \\
}

\begin{document}

\maketitle

% ---------------------------------------------------------------
\begin{abstract}
Food image classification is a challenging computer vision problem due to high
intra-class visual variance and strong inter-class similarity across dish
categories. In this work, we apply transfer learning with an EfficientNet-B0
backbone fine-tuned on the Food-101 dataset to build a robust 101-class food
image classifier. Our proposed model achieved an accuracy of \textbf{65.72\%} on
the test set, outperforming a baseline CNN trained from scratch by \textbf{3.65\%}.
This demonstrates the significant benefit of ImageNet pre-training for food recognition tasks.
\end{abstract}

% ---------------------------------------------------------------
\section{Introduction}

\paragraph{Motivation}
Automated food recognition has direct practical applications in dietary tracking,
nutritional analysis, and smart restaurant checkout systems. Manual food logging
is both inaccurate and time-consuming; a reliable image classifier eliminates
this friction and enables scalable health monitoring pipelines~\citep{bossard2014food101}.
With over 1 billion smartphone users capable of capturing food images, the need
for accurate, fast food classification is more pressing than ever.

\paragraph{Problem Definition}
We formulate food recognition as a supervised multi-class image classification
task. Given an input RGB image $\mathbf{x} \in \mathbb{R}^{H \times W \times 3}$,
the model predicts a class label $\hat{y} \in \{1, 2, \ldots, 101\}$
corresponding to one of 101 food categories in the Food-101 dataset. The task
requires robust feature extraction to handle significant variations in plating,
lighting, camera angles, and food presentation.

\paragraph{Contributions}
The main contributions of this project are:
\begin{itemize}
    \item We apply transfer learning with a pre-trained EfficientNet-B0 backbone
          to Food-101, achieving 65.72\% classification accuracy on a
          101-class benchmark with frozen backbone fine-tuning.
    \item We compare EfficientNet-B0 against a baseline CNN trained from scratch,
          demonstrating a 3.65\% improvement through ImageNet pre-training.
    \item We conduct a detailed per-class error analysis using confusion matrices
          to identify the most challenging food category pairs and provide insights
          into visual similarities that confound the classifiers.
\end{itemize}

% ---------------------------------------------------------------
\section{Method}

\subsection{Algorithm Description}

\paragraph{Baseline CNN}
The baseline model consists of a ResNet18 architecture~\citep{he2016deep} trained from scratch on Food-101.
ResNet18 is a lightweight residual network with 18 layers that learns hierarchical
feature representations through skip connections. The model is trained to minimize cross-entropy loss:
\[
    \mathcal{L} = -\frac{1}{N}\sum_{i=1}^{N} \log \frac{e^{z_{y_i}}}{\sum_{j=1}^{101} e^{z_j}}
\]
where $z_j$ is the logit for class $j$ and $y_i$ is the ground-truth label. The baseline
provides a controlled comparison point to measure the benefit of pre-training.

\paragraph{EfficientNet-B0 with Transfer Learning}
EfficientNet scales network depth, width, and resolution uniformly using a
compound coefficient~\citep{tan2019efficientnet}. We use a pre-trained
EfficientNet-B0 backbone with ImageNet weights, freeze all convolutional layers,
and replace the final classification head with a linear layer of output
dimension 101. By freezing the backbone, we leverage pre-learned features (learned from ImageNet~\citep{krizhevsky2012imagenet}) while
allowing the classification head to adapt to the food recognition task. The model
is trained by minimizing the same cross-entropy loss as the baseline, applied only
to the trainable classification head parameters.

\subsection{Overall Framework}

Figure~\ref{fig:pipeline} illustrates the end-to-end pipeline, which consists
of four stages: (1) data preprocessing and augmentation, (2) feature extraction via the
frozen EfficientNet-B0 backbone pre-trained on ImageNet, (3) classification head training
on Food-101 labels, and (4) evaluation of test set predictions.

\begin{figure}[ht]
    \centering
    \includegraphics[width=0.95\linewidth]{framework.png}
    \caption{Overall framework for Food-101 image classification pipeline. Raw images are
    resized and normalized before being passed through a frozen EfficientNet-B0 backbone.
    The classification head maps extracted features to one of 101 food categories.}
    \label{fig:pipeline}
\end{figure}

\subsection{Design Choices}

All images are resized to $224 \times 224$ pixels to match the input resolution
expected by EfficientNet-B0 and to ensure consistent memory usage during
batched training. Pixel values are normalized using the ImageNet statistics
$\mu = (0.485, 0.456, 0.406)$ and $\sigma = (0.229, 0.224, 0.225)$
to align the input distribution with the pre-training regime. During training,
we apply random horizontal flips (50\% probability), random crops, and colour jitter
to improve generalization and reduce overfitting on domain-specific visual patterns~\citep{shorten2019augmentation}.
These augmentations are applied only to the training set; the validation and test
sets use centre crops without augmentation to ensure fair evaluation.

% ---------------------------------------------------------------
\section{Experiments}

\subsection{Experimental Setup}

The Food-101 dataset~\citep{bossard2014food101} contains 101,000 images across
101 food categories, with approximately 1,000 images per class at variable resolution.
We used the official train/test split of 75,750 training images and 25,250 test images.
An additional 10\% of the training set (approximately 7,500 images) was reserved as a
validation set to tune hyperparameters and select checkpoints. All experiments were conducted
on a single NVIDIA T4 GPU (16 GB VRAM) via Google Colab using PyTorch 1.10 and torchvision 0.11.

Key hyperparameters for both models are listed in Table~\ref{tab:hyperparams}.
Both models were optimized using the Adam optimizer with the same learning rate and batch
size constraints to ensure fair comparison. The baseline CNN was trained for 50 epochs,
while EfficientNet-B0 (with frozen backbone) converged faster and was trained for 30 epochs.

\begin{table}[h]
    \caption{Training hyperparameters for both models on Food-101.}
    \label{tab:hyperparams}
    \centering
    \begin{tabular}{lcc}
        \toprule
        Hyperparameter & Baseline CNN & EfficientNet-B0 \\
        \midrule
        Optimizer      & Adam         & Adam \\
        Learning Rate  & 1.0e-3       & 1.0e-3 \\
        Batch Size     & 64           & 32 \\
        Epochs         & 50           & 30 \\
        Image Size     & 224x224      & 224x224 \\
        Weight Decay   & 1.0e-5       & 1.0e-5 \\
        \bottomrule
    \end{tabular}
\end{table}

\subsection{Model Comparison Results}

Table~\ref{tab:results} presents the performance of the baseline CNN and our
EfficientNet-B0 model on the Food-101 test set across four metrics.

\begin{table}[h]
    \caption{Model comparison on the Food-101 test set. EfficientNet-B0 with transfer
    learning outperforms the baseline CNN across all metrics, demonstrating the 
    substantial benefit of ImageNet pre-training for food image classification.}
    \label{tab:results}
    \centering
    \begin{tabular}{lcccc}
        \toprule
        Model           & Accuracy (\%) & Precision & Recall & F1-Score \\
        \midrule
        Baseline CNN    & 62.07         & 0.6217    & 0.6207 & 0.6190 \\
        EfficientNet-B0 & 65.72         & 0.6534    & 0.6572 & 0.6533 \\
        \midrule
        \textbf{Improvement} & \textbf{+3.65} & \textbf{+0.0317} & \textbf{+0.0365} & \textbf{+0.0343} \\
        \bottomrule
    \end{tabular}
\end{table}

\subsection{Justification and Analysis}

EfficientNet-B0 achieves significantly higher accuracy than the baseline CNN
(65.72\% vs 62.07\%, a 3.65 percentage point improvement). This improvement
is consistent across precision, recall, and F1-score metrics, all favoring the
transfer learning approach. The consistent gains across all metrics indicate that
EfficientNet-B0 provides both better overall accuracy and more balanced per-class
performance.

The superior performance of EfficientNet-B0 can be attributed to its ImageNet-pretrained
feature representations, which capture generic low-level visual patterns (edges, textures)
and high-level semantic structures (object parts, shapes) that generalize well to food images.
In contrast, the baseline CNN must learn these patterns from Food-101 alone, which is
limited by the training budget and domain specificity of food images. The frozen backbone
approach provides an effective regularization mechanism, preventing overfitting while
allowing the shallow classification head to adapt to the 101-class food recognition task.

Figure~\ref{fig:confusion} shows the confusion matrix for EfficientNet-B0 on the test set.
Analysis of the top misclassifications reveals that the most common errors occur between
visually similar food categories. For instance, certain plated dishes (e.g., beef carpaccio
and ceviche) can be easily confused due to similar presentation and colour palettes.
Additionally, categorical confusion is observed within ``batter-based'' foods and within
``salad'' varieties, suggesting that the model struggles with fine-grained distinctions
within visually cohesive food groups. These errors reflect inherent ambiguities in food
photography and plating conventions rather than fundamental limitations of the architecture.

\begin{figure}[ht]
    \centering
    \includegraphics[width=0.95\linewidth]{confusion_matrix.png}
    \caption{Confusion matrix for EfficientNet-B0 on the Food-101 test set showing
    classification accuracy across all 101 food categories. Darker regions indicate
    more frequent misclassifications and highlight visually similar food pairs that
    are challenging for the model to distinguish.}
    \label{fig:confusion}
\end{figure}

The trade-off between the two approaches is clear: the baseline CNN requires longer
training time and greater computational tuning but can in principle achieve higher accuracy
with sufficient hyperparameter optimization and data augmentation. EfficientNet-B0, while
achieving better performance with frozen weights, would benefit from selective layer unfreezing
in deeper backbone layers to allow adaptation to food-specific features. However, given the
constraints of a single T4 GPU and the time budget for this project, the frozen backbone
proved an effective and practical choice.

% ---------------------------------------------------------------
\section{Conclusion}

\paragraph{Summary}
We presented a transfer learning pipeline for food image classification on Food-101,
comparing a baseline ResNet18 trained from scratch against a pre-trained EfficientNet-B0
with frozen backbone. Our EfficientNet-B0 model achieved a test accuracy of \textbf{65.72\%}
and an F1-score of \textbf{0.6533}, outperforming the baseline CNN by \textbf{3.65\%} in
accuracy. The analysis demonstrates that leveraging ImageNet pre-training yields consistent
improvements across all evaluation metrics.

\paragraph{Limitations and Future Work}
A primary limitation of our approach is that the EfficientNet backbone is fully frozen during
training, which may prevent the model from learning food-specific visual features that diverge
from generic ImageNet patterns. Given more computational time and resources, we would explore
progressive unfreezing of deeper layers (e.g., unfreezing the final 2--3 residual blocks) to
improve adaptation while maintaining the regularization benefits of pre-training. Additionally,
experimenting with larger EfficientNet variants (e.g., EfficientNet-B1 through B3) could
further improve performance, though at increased computational cost. Finally, ensemble methods
combining multiple models and test-time augmentation could provide modest accuracy gains at
modest computational overhead.

% ---------------------------------------------------------------
\section*{Distribution of Work}
\begin{itemize}
    \item \textbf{Jasia Azreen:} Data pipeline design, preprocessing and augmentation
          implementation, baseline ResNet18 model implementation and training.
    \item \textbf{Sofiya Spolitak:} EfficientNet-B0 transfer learning architecture setup,
          training loop implementation, and hyperparameter tuning.
    \item \textbf{John Song:} Evaluation metrics computation, confusion matrix
          generation, error analysis, pipeline diagram design, and final report assembly.
\end{itemize}

% ---------------------------------------------------------------
\bibliographystyle{plainnat}
\bibliography{references}

\end{document}
